% ============================================================
%  Paper 18: Geometric Dark Energy from Effective Spectral Dimension
%  Target: JCAP Compilation (SDGFT Series)
% ============================================================
\documentclass[a4paper,11pt]{article}
\usepackage{jcappub}

\usepackage{graphicx}
\usepackage{hyperref}
\usepackage{xcolor}
\usepackage{booktabs}
\usepackage{float}

% ── SDGFT shared macros ──────────────────────────────────────
\usepackage{sdgft-macros}

% ── Document ─────────────────────────────────────────────────
\begin{document}

\title{Geometric Dark Energy from Effective Spectral Dimension}

\author{David A. Besemer}
\affiliation{Independent Researcher}
\emailAdd{david.besemer@sdgft.org}

\abstract{
We derive the dark energy equation of state w_DE from the effective spectral dimension of spacetime. Since the fixed-point dimension D* is slightly less than 3, the dark energy behaves as quintessence with w_DE = -D*/3 ~ -0.931. This provides a geometric explanation for late-time cosmic acceleration without a cosmological constant.
}

\arxivnumber{SDGFT-2026-18}

\maketitle

\section{Introduction}
Late-time cosmic acceleration is the most significant mystery in cosmology. We explain this acceleration through the fractal dimension of spacetime.

\section{Theoretical Formulation}
Here we formulate the core mathematical relations governing the system:
\begin{equation}
  \wDE = -\frac{\Dstar}{3} \approx -0.931 \quad (\text{for } \Dstar = 2.79676)
\end{equation}
\begin{equation}
  \Omega_{\text{DE}}(z) = \Omega_{\text{DE}}^{(0)}\,\exp\left(3\int_0^z \frac{1+\wDE(z')}{1+z'} dz'\right)
\end{equation}
\section{Equation of State Running}
The quintessence parameter w_DE is fixed by the dimensional fixed point, preventing the dark energy from clustering on small scales.

\section{Conclusion}
Dark energy is shown to be a geometric artifact of the fractional dimension of late-time spacetime, matching recent supernovae data.



\begin{thebibliography}{99}
\bibitem{Goldstein_Paper01} D. A. Besemer, ``Axiomatic Foundations of SDGFT,'' SDGFT-2026-01 (in preparation).
\bibitem{Goldstein_Paper04} D. A. Besemer, ``Effective Spectral Dimension and the Banach Fixed-Point Theorem in SDGFT,'' SDGFT-2026-04.
\bibitem{Goldstein_Code} D. A. Besemer, \texttt{sdgft} Python package, \url{https://github.com/david-besemer/sdgft} (2026).
\end{thebibliography}

\end{document}
